\input{../tex-inputs/latex-leseansicht-vorspann}

\section[Arthur und Olga Schnitzler an Gustav Schwarzkopf, 10. 8. 1907]{L04471 Arthur und Olga Schnitzler an Gustav Schwarzkopf, 10. 8. 1907}
\nopagebreak\mylabel{L04471v}
\rehead{ }\normalsize\beginnumbering\briefempfaengerindex{Schwarzkopf, Gustav@\textsc{Schwarzkopf, Gustav}!zzzSchnitzler, Olga@\emph{von Olga Schnitzler}!1907-08-101@{10. 8. 1907}|(be}\briefempfaengerindex{Schwarzkopf, Gustav@\textsc{Schwarzkopf, Gustav}!zzzSchnitzler, Arthur@\emph{von Arthur Schnitzler}!1907-08-101@{10. 8. 1907}|(be}
\toendnotes[C]{\smallbreak\pagebreak[2]}
\correspDesc{Versand  durch Arthur Schnitzler, Olga Schnitzler am 10. 8. 1907 in Wildbad Waldbrunn
\newline{}Übermittlung  in Welsberg-Taisten
\newline{}Erhalt  durch Gustav Schwarzkopf im Zeitraum [11. 8. 1907
                  – 15. 8. 1907?] in Wien}\toendnotes[C]{\smallbreak}
\Standort{DLA, A:Schnitzler, HS.1985.1.1897.}
\physDesc{Bildpostkarte, 215 Zeichen
\newline{}Handschrift Arthur Schnitzler: Bleistift, deutsche Kurrent
\newline{}Handschrift Olga Schnitzler: Bleistift, lateinische Kurrent
\newline{}Versand: Stempel: »\nobreak{}\oindex{Welsberg-Taisten@\textbf{Welsberg-Taisten}|pwk}\textcolor{gray}{W}els\textcolor{gray}{be}r\textcolor{gray}{g}\nobreak{}«.  }\toendnotes[C]{\smallbreak}\pstart{}{\pb}\textcolor{gray}{\textbf{An}} Hrn Guſtav Schwarzkopf\pend{}\pstart{}Wien I\oindex{I., Innere Stadt@\textbf{I., Innere Stadt}|pw}\pend{}\pstart{}I. Tiefer Graben 23\oindex{Wien@\textbf{Wien}!I., Innere Stadt@\textbf{I., Innere Stadt}!Tiefer Graben 23@\textbf{Tiefer Graben 23}|pw}\pend{}{\bigskip}
\pstart
           \noindent{}\centering{}{\pb}\textcolor{gray}{\textbf{Hôtel und Pension Wildbad Waldbrunn b. Welsberg
                     (Pustertal)}}\oindex{Wildbad Waldbrunn@\textbf{Wildbad Waldbrunn}|pw}\pend
           \vspace{1em}
\pstart
           \noindent{}{\pb}\label{T_L04471-1v}\edtext{10. 8. 907, noch i{\geminationm}er, lieber
                  Guſtav{ }ſind}{\lemma{\textnormal{\emph{10. 8. 907, … sind}}}\Cendnote{\textnormal{entlang des linken
                  Randes geschrieben}}}\label{T_L04471-1}{ }\label{T_L04471-2v}\edtext{wir oben}{\lemma{\textnormal{\emph{wir oben}}}\Cendnote{\textnormal{ab hier am rechten Rand}}}\label{T_L04471-2};{ }ſeit 8 Tagen iſt Mama\pwindex{Schnitzler, Louise 8.\,7.\,1840 Kőszeg – 9.\,9.\,1911 Wien@\textsc{Schnitzler, Louise} (8.\,7.\,1840 Kőszeg – 9.\,9.\,1911 Wien)|pwv} bei uns. Wir wollten
               vielleicht nach  20.{ }ſüdwärts, Meran\oindex{Meran@\textbf{Meran}|pw}, od Riva\oindex{Riva del Garda@\textbf{Riva del Garda}|pw}\pend
           
\pstart
           herzlichſt{\\[\baselineskip]}Ihr \spacefill\mbox{A.}{\\[\baselineskip]}{[}hs. O. Schnitzler:{]} \label{T_L04471-3v}\edtext{Herzliche Grüsse! \spacefill\mbox{Olga}}{\lemma{\textnormal{\emph{Herzliche Grüsse! Olga}}}\Cendnote{\textnormal{Olga Schnitzlers\pwindex{Schnitzler, Olga 17.\,1.\,1882 Wien – 13.\,1.\,1970 Lugano@\textsc{Schnitzler, Olga} (17.\,1.\,1882 Wien – 13.\,1.\,1970 Lugano)|pwk} Gruß steht am unteren Rand
                  unter dem Bild.}}}\label{T_L04471-3}\pend
           \leftskip=0em{}\selectlanguage{ngerman}\endnumbering\briefempfaengerindex{Schwarzkopf, Gustav@\textsc{Schwarzkopf, Gustav}!zzzSchnitzler, Olga@\emph{von Olga Schnitzler}!1907-08-101@{10. 8. 1907}|)be}\briefempfaengerindex{Schwarzkopf, Gustav@\textsc{Schwarzkopf, Gustav}!zzzSchnitzler, Arthur@\emph{von Arthur Schnitzler}!1907-08-101@{10. 8. 1907}|)be}\mylabel{L04471h}
\begin{anhang}
\end{anhang}\newcommand{\dateiname}{L04471}\newcommand{\titel}{Arthur und Olga Schnitzler an Gustav Schwarzkopf, 10. 8. 1907}\newcommand{\editorInnen}{Herausgegeben von Jahnke, SelmaMüller, Martin Anton}\input{../tex-inputs/latex-leseansicht-abspann}
